\chapter{JADWAL PENELITIAN}

Penelitian ini direncanakan untuk dilaksanakan selama delapan bulan, dari bulan Januari 2026 hingga Agustus 2026. Jadwal penelitian disusun berdasarkan tahapan studi kuasi-eksperimen: persiapan lingkungan, kalibrasi baseline, pelaksanaan eksperimen, analisis statistik, dan penulisan laporan.

Tabel \ref{tab:jadwal} menunjukkan rencana jadwal pelaksanaan penelitian berbasis bulan. Simbol $\bullet$ menunjukkan bahwa kegiatan dilaksanakan pada bulan tersebut.

\begin{table}[H]
  \centering
  \caption{Jadwal Penelitian (Gantt Chart)}
  \label{tab:jadwal}
  \begin{tabular}{|l|c|c|c|c|c|c|c|c|}
    \hline
    \multirow{2}{*}{\textbf{Kegiatan}} & \multicolumn{8}{c|}{\textbf{Bulan}} \\
    \cline{2-9}
     & \textbf{Jan} & \textbf{Feb} & \textbf{Mar} & \textbf{Apr} & \textbf{Mei} & \textbf{Jun} & \textbf{Jul} & \textbf{Agu} \\
    \hline
    Studi Literatur & $\bullet$ & $\bullet$ & & & & & & \\
    \hline
    Perancangan Lingkungan & & $\bullet$ & $\bullet$ & & & & & \\
    \hline
    Implementasi ArgoCD + Baseline Manual & & & $\bullet$ & $\bullet$ & & & & \\
    \hline
    Persiapan Baseline \& Kalibrasi & & & & $\bullet$ & $\bullet$ & & & \\
    \hline
    Eksperimen \& Pengumpulan Data (E1--E4) & & & & & $\bullet$ & $\bullet$ & & \\
    \hline
    Analisis Statistik \& Interpretasi & & & & & & $\bullet$ & $\bullet$ & \\
    \hline
    Penulisan Laporan & $\bullet$ & $\bullet$ & $\bullet$ & $\bullet$ & $\bullet$ & $\bullet$ & $\bullet$ & $\bullet$ \\
    \hline
    Seminar / Sidang & & & & & & & & $\bullet$ \\
    \hline
  \end{tabular}
\end{table}

\section{Penjelasan Kegiatan}

\begin{enumerate}
    \item \textbf{Studi Literatur (Januari -- Februari 2026)}
    
    Kegiatan ini mencakup pengumpulan dan analisis literatur mengenai: \textit{configuration drift} pada Kubernetes, evaluasi empiris GitOps, metode kuasi-eksperimen pada riset sistem, dan arsitektur InvenioRDM sebagai beban uji. Sumber literatur meliputi jurnal ilmiah (IEEE, ACM, Springer), dokumentasi resmi, buku teks, dan artikel teknis. Tujuan dari fase ini adalah membangun landasan teoretis dan menetapkan hipotesis yang dapat diuji.
    
    \item \textbf{Perancangan Lingkungan (Februari -- Maret 2026)}
    
    Tahap ini meliputi perancangan arsitektur Lingkungan A (deployment manual/\texttt{kubectl}) dan Lingkungan B (deployment GitOps/ArgoCD) pada klaster Kubernetes yang sama. Perancangan meliputi: definisi variabel independen dan dependen, pemilihan tujuh skenario drift, rancangan tabel rubrik traceability, dan spesifikasi protokol pengukuran.
    
    \item \textbf{Implementasi ArgoCD + Baseline Manual (Maret -- Mei 2026)}
    
    Tahap implementasi meliputi: konfigurasi ArgoCD pada klaster, pembuatan repositori GitOps dengan manifest deklaratif, konfigurasi baseline manual (alur kerja \texttt{kubectl}), implementasi skenario drift (script injeksi otomatis), serta konfigurasi logging dan monitoring (Prometheus, Grafana, Kubernetes audit logs, ArgoCD notifications). Pada akhir fase ini, kedua lingkungan siap untuk eksperimen.
    
    \item \textbf{Persiapan Baseline \& Kalibrasi (April -- Mei 2026)}
    
    Sebelum eksperimen dimulai, fase kalibrasi dilakukan untuk memastikan validitas internal. Kegiatan meliputi: (a)~deployment InvenioRDM pada kedua lingkungan; (b)~verifikasi baseline bersih; (c)~pengukuran awal konsistensi; (d)~verifikasi penggunaan sumber daya node di bawah 70\%; dan (e)~uji coba awal untuk memvalidasi protokol pengukuran.
    
    \item \textbf{Eksperimen dan Pengumpulan Data (Mei -- Juni 2026)}
    
    Pelaksanaan empat eksperimen sesuai desain pada Bab III. E1: konsistensi deployment (n = 5 per skenario). E2: waktu deteksi dan pemulihan drift (n = 10 per skenario). E3: jumlah intervensi manual (n = 5 per skenario). E4: traceability (n = 5 per skenario, lengkap dengan rubrik dan MTTI). Data dikumpulkan dalam bentuk \textit{timestamp}, status aplikasi, dan log. Skenario diacak untuk setiap lingkungan untuk memitigasi \textit{order effect}.
    
    \item \textbf{Analisis Statistik dan Interpretasi (Juni -- Juli 2026)}
    
    Analisis data meliputi: uji normalitas (Shapiro--Wilk), uji hipotesis (paired t-test atau Wilcoxon signed-rank), perhitungan ukuran efek (Cohen's d atau Cliff's delta), dan interpretasi hasil dalam konteks H$_0$ dan H$_1$. Pembahasan mencakup identifikasi skenario drift yang paling berdampak, faktor yang mempengaruhi hasil, serta validitas internal dan eksternal penelitian.
    
    \item \textbf{Penulisan Laporan (Januari -- Agustus 2026)}
    
    Penulisan laporan dilakukan secara paralel sepanjang penelitian. Setiap tahap menghasilkan bab atau bagian yang sesuai. Draft proposal diselesaikan pada pertengahan semester, draft skripsi lengkap pada akhir penelitian.
    
    \item \textbf{Seminar / Sidang (Agustus 2026)}
    
    Presentasi dan sidang tugas akhir di depan dosen penguji, mencakup paparan hasil eksperimen, analisis statistik, dan kesimpulan.
\end{enumerate}
